% !TEX root = ../main.tex
\chapter{Architettura di Lean 4}
\label{chap:architecture}

% TODO: add cappello

\section{Pipeline di compilazione}
L'architettura di Lean 4 si distingue nel panorama dei linguaggi moderni per la sua duplice natura, contemporaneamente un assistente alla dimostrazione interattivo 
(\emph{interactive theorem prover}) e un linguaggio di programmazione general-purpose. Per supportare questa ambivalenza, il processo di compilazione non è un semplice 
traduttore da codice sorgente a linguaggio macchina, ma una sofisticata pipeline che accompagna il codice attraverso molteplici livelli di astrazione. La pipeline di 
compilazione, dal file di testo fino all'esecuzione, può essere compsreso analizzando tre macro-fasi distinte, ognuna con responsabilità rigorosamente separate. 
\begin{description}[leftmargin=1.5em, style=nextline, font=\bfseries]
    \item [Frontend]
        Il responsabile dell'interazione con l'utente. Qui il testo viene analizzato, le eventuali macro vengono espanse (risolvendo lo ``zucchero sintattico'') 
        e il codice viene elaborato fino a diventare un'espressione formale pura. È la fase che colma il divario tra l'intuizione umana e il rigore meccanico.
    \item [Kernel]
        Il responsabile della corretezza matematica del sistema, agisce come un controllore inflessibile dove la sua unica funzione è verificare la logica delle espressioni 
        prodotte dal frontend. Costituisce la \textit{Trusted Computing Base} (TCB) di Lean, ovvero l'unica parte del sistema di cui ci si deve fidare per garantire la validità 
        delle dimostrazioni. Per garantire una buon livello di affidabilità, il kernel è progettato per essere il più piccolo e semplice possibile.
    \item [Backend]
        Una volta che il kernel ha validato la correttezza logica, il backend si occupa di ottimizzare e compilare il codice. Questa fase è completamente separata dalla logica 
        formale. Tale processo avviene esclusivamente per le espressioni computazionali, eliminando ogni informazione matematica superflua e generando codice macchina ``efficiente''.
\end{description}
Questa netta separazione delle preoccupazioni (\textit{separation of concerns}) garantisce che l'efficienza dell'esecuzione non comprometta mai l'assoluta affidabilità 
delle dimostrazioni matematiche.

\subsection{Frontend dal parsing alla laborazione}
Il \textit{Frontend} di Lean 4 è l'interfaccia architetturale che media tra l'utente (tramite l'editor) e il rigoroso mondo del Kernel. A differenza dei compilatori 
tradizionali, il frontend di Lean è progettato per essere interattivo, incrementale ed estensibile a runtime dall'utente stesso. Esso trasforma un flusso di caratteri 
testuali in un'espressione validabile (\texttt{Expr}), attraversando tre fasi computazionali distinte: parsing, espansione delle macro ed elaborazione.

\subsubsection{Lexing, Parsing e l'Albero Sintattico Concreto (CST)}
Il processo ha inizio con il \emph{Parser}, il quale analizza il codice sorgente. A differenza di molti linguaggi funzionali che utilizzano generatori di parser statici 
(come Yacc o Bison), Lean 4 implementa un parser a discesa ricorsiva estensibile, basato su un algoritmo noto come \textit{Pratt Parsing} (parsing a precedenza di operatori). 
Questo permette all'utente di definire nuove notazioni sintattiche al volo, assegnando loro specifici livelli di precedenza. Il risultato del parsing non è un tradizionale
\textit{Abstract Syntax Tree} (AST), ma un \textit{Albero Sintattico Concreto} (CST), denotato in Lean dal tipo induttivo \texttt{Syntax}. Il tipo \texttt{Syntax} è progettato 
per essere una rappresentazione completamente fedele (fully faithful) del file sorgente. Nello specifico, esso memorizza simultaneamente la struttura gerarchica dei nodi sintattici, 
i metadati lessicali (\texttt{SourceInfo}) per tracciare le esatte coordinate di riga e colonna, e le informazioni tipografiche, preservando gli spazi bianchi (\textit{whitespaces}) 
e i commenti adiacenti a ogni token. Questa scelta di conservare le informazioni non semantiche è cruciale poiche permette al server linguistico di Lean (LSP) di fornire feedback 
visivo, autocompletamento e \textit{refactoring} testuale direttamente nell'editor, garantendo che la manipolazione automatica del codice non ne distrugga la formattazione originale.

\subsubsection{Espansione delle macro}
Una volta generato l'albero concreto tramite il tipo induttivo \texttt{Syntax}, il controllo passa al sistema delle Macro. In Lean 4, una macro è una funzione
pura che mappa sintassi in sintassi (\texttt{Syntax $\rightarrow$ MacroM Syntax}), operando all'interno della monade \texttt{MacroM}. L'obiettivo delle macro è 
la "dezuccherizzazione" (\textit{desugaring}), costrutti di alto livello o notazioni create dall'utente vengono meccanicamente tradotte in costrutti più primitivi. 
Questo processo avviene prima che il compilatore tenti di risolvere i tipi del programma. Ad esempio, un costrutto come la notazione per le liste \texttt{[a, b, c]} 
non esiste nel nucleo del linguaggio. Tale notazione viene espansa tramite macro in funzioni ricorsive o catene di costruttori (\texttt{List.cons}). Questo approccio 
modulare mantiene l'elaboratore semantico snello, sollevandolo dal dover conoscere ogni singola notazione derivata.

% \subsubsection{Elaboratore}
% TODO

\subsection{Kernel e Trusted Computing Base (TCB)}
L'elaborazione descritta in precedenza culmina con la generazione di un termine \texttt{Expr} completamente esplicito. A questo punto, il termine lascia il dominio del frontend 
e viene sottoposto al giudizio del \textbf{Kernel}. In Lean 4, il kernel rappresenta la \textit{Trusted Computing Base} (TCB) del sistema, ovvero l'unica, vitale componente software 
della cui correttezza l'utente deve fidarsi per garantire la solidità matematica delle dimostrazioni.

\subsubsection{Principio di de Bruijn}
Questa netta separazione architetturale incarna il \emph{Principio di de Bruijn}. Tale principio fondazionale per gli assistenti alla dimostrazione stabilisce che un sistema di verifica 
formale debba separare rigorosamente la fase di costruzione della prova dalla fase di verifica della stessa. Il frontend (comprendente parser, macro e algoritmi di elaborazione) è costituito 
da decine di migliaia di righe di codice complesso. Per la natura intrinseca dello sviluppo software, esso è statisticamente incline a contenere bug informatici e viene pertanto considerato 
non fidato (\textit{untrusted}). Non possiede alcuna autorità per dichiarare un teorema valido: si limita a fungere da intermediario, proponendo una potenziale prova.

\subsubsection{Verifica meccanica}
Per mitigare il rischio di errori logiche, il kernel è progettato all'insegna del minimalismo e dell'isolamento. Attualmente implementato in C++ (per ragioni storiche e di performance), il kernel 
è un pezzo di codice volutamenete compatto. Il suo unico compito è l'esecuzione di un algoritmo deterministico di \textit{Type Checking}. A differenza dell'elaboratore, il kernel è ``cieco'' alle 
euristiche, non effettua alcuna inferenza, non sa cosa siano le typeclass e metavariabili non risolte. Esso si limita a ispezionare il termine \texttt{Expr} fornito, verificando meccanicamente e 
passo dopo passo che esso rispetti gli assiomi e le regole deduttive del \textit{Calculus of Inductive Constructions} (CIC), discussi del capitolo \ref{chap:cic}. Se il type checker convalida il 
termine, la dichiarazione viene accettata e memorizzata nell'ambiente (\textit{Environment}) come matematicamente ineccepibile. Qualora la \texttt{Expr} contenga un errore o derivi da un bug del 
frontend, il kernel la rifiuterà inflessibilmente, impedendo a priori la corruzione del sistema formale.

\subsection{Backend dal'ottimizzazione all'esecuzione}
Per le espressioni dotate di semantica operazionale (come le funzioni introdotte dalle parole chiave \texttt{def} o \texttt{abbrev}), la validità logica stabilita dal kernel rappresenta solo un traguardo 
intermedio. Per poter permettere a Lean 4 di competere in termini di prestazioni con i linguaggi di programmazione industriali, il codice funzionale deve essere compilato in istruzioni macchina altamente 
ottimizzate. Tale compito è l'esclusiva responsabilità del backend. Il passaggio dal dominio logico a quello computazionale avviene attraverso una pipeline di \emph{lowering} che degrada progressivamente 
l'astrazione matematica per modellare i rigidi vincoli dell'hardware sottostante.

\subsubsection{Eliminazione del contenuto proposizionale}
Il termine \texttt{Expr} approvato dal kernel è completamente annotato di dettagli logici, come tipi dipendenti complessi e intere dimostrazioni di correttezza intrise nel codice. Tali dettagli sebbene cruciali 
per la coretteza del sistema formale a \textit{runtime} non sono necessari per l'esecuzione effettiva del programma. Pertanto, la prima operazione del backend è la cancellazione dei tipi (\texttt{Type Erasure}). 
Il compilatore analizza l'espressione e sopprime tutti i sotto-termini appartenenti all'universo proposizionale (\texttt{Prop}). Poiché le dimostrazioni non hanno alcun effetto sul comportamento a runtime, esse 
vengono rimosse e sostituite internamente con un valore fittizio e unitario (denotato come \texttt{lcErased} o \scalebox{0.8}{$\blacksquare$}). Questo processo garantisce che il codice generato non paghi alcun 
\textit{overhead} prestazionale per il rigore logico preteso nelle fasi precedenti.

\subsubsection{LCNF, A-Normal Form per l'ottimizzazione funzionale}

\subsubsection{Lean IR, Gestione della memoria e in-place mutation}

\subsubsection{Compilazione nativa mediante delega architetturale}
L'ultima parte della pipeline è sorprendentemente pragmatica. La Lean IR è progettata per essere semanticamente vicina al linguaggi imperativi, in particolare al linguaggio C. Il backend mappa i blocchi base 
e le istruzioni di \textit{reference counting} in codice sorgente C standard, appoggiandosi a un \textit{runtime} minimale scritto in C++. Tale runtime fornisce i servizi di base per la gestione della memoria, 
come l'allocazione dinamica e la deallocazione (tramite la struttura \texttt{lean\_object}), interi a precisione arbitraria e gestione della processazione parallela. Al fine mi massimizzare la portabilità e 
minimizzare la complessità, il backend non tenta di generare codice macchina direttamente. Tale compito è delegato a compilatori industriali consolidati (come GCC o LLVM/Clang). Questi strumenti prendono il 
codice C generato e vi applicano decenni di ottimizzazioni hardware di basso livello, producendo infine eseguibili binari nativi ad altissime prestazioni.